\chapter{bipartite}
\label{ch:bipartite}

\begin{Proposition}
\label{prop:bip-st-ind-hard}
    \(stable\)-Unabhängigkeit ist auch für bipartite AFs
    \(\Pi_2^{\mathrm P}\)-vollständig.
\end{Proposition}

\begin{proof}
Die Zugehörigkeit zu \(\Pi_2^{\mathrm P}\) gilt bereits für beliebige AFs.
Tatsächlich kann Nicht-Unabhängigkeit durch zwei stabile Labellings
\(\lambda_1,\lambda_2\) bezeugt werden, für die kein stabiles Labelling
\(\lambda\) die geforderten Einschränkungen kombiniert. Dies lässt sich
durch eine Formel der Form
\(\exists\lambda_1\exists\lambda_2\forall\lambda\) mit polynomiell
prüfbarem Prädikat ausdrücken. Nicht-Unabhängigkeit liegt daher in
\(\Sigma_2^{\mathrm P}\) und \(stable\)-Unabhängigkeit folglich in
\(\Pi_2^{\mathrm P}\).

Für die Härte reduzieren wir vom \(\Pi_2^{\mathrm P}\)-vollständigen Problem
\textsc{True \(\forall\exists\)-3-CNF}. Sei
\[
    \Psi=\forall X\exists Y\,\varphi
\]
eine Instanz dieses Problems.

Zunächst bringen wir \(\varphi\) in eine Form, in der jede Klausel entweder
ausschließlich positive oder ausschließlich negative Literale enthält.
Sei \(C\) eine gemischte Klausel. Wir schreiben
\[
    C=P_C\lor N_C,
\]
wobei \(P_C\) die Disjunktion der positiven und \(N_C\) die Disjunktion
der negativen Literale von \(C\) bezeichnet. Für \(C\) führen wir eine
frische existentielle Variable \(z_C\) ein und ersetzen \(C\) durch
\[
    (P_C\lor z_C)\land(N_C\lor\neg z_C).
\]
Für jede Belegung der ursprünglichen Variablen gilt
\[
    C
    \quad\Longleftrightarrow\quad
    \exists z_C\,
    \bigl((P_C\lor z_C)\land(N_C\lor\neg z_C)\bigr).
\]
Ist \(P_C\) wahr, kann \(z_C=0\) gewählt werden. Ist \(N_C\) wahr, kann
\(z_C=1\) gewählt werden. Sind beide Teile falsch, erzwingt die erste
Klausel \(z_C=1\), während die zweite \(z_C=0\) erzwingt.

Jede gemischte Klausel erhält eine eigene neue Variable, und alle neuen
Variablen werden dem existentiellen Block hinzugefügt. Die erhaltene QBF
ist daher genau dann gültig wie \(\Psi\). Da die ursprünglichen Klauseln
höchstens drei Literale enthalten, haben auch die neuen Klauseln höchstens
drei Literale. Die Transformation ist polynomial. Im Folgenden bezeichnen
wir die transformierte QBF wieder mit
\[
    \Psi=\forall X\exists Y\,\varphi.
\]
Die Klauselmenge von \(\varphi\) zerfällt nun in die Menge
\(\mathcal C^+\) der positiven und die Menge \(\mathcal C^-\) der
negativen Klauseln.

Wir führen zwei neue Variablen \(p,n\) ein und setzen
\[
    V=X\cup Y\cup\{p,n\}.
\]
Das AF
\[
    F_\Psi=(\mathcal A_\Psi,R_\Psi)
\]
ist durch
\[
\begin{aligned}
    \mathcal A_\Psi
    ={}&
    \{v,\overline v\mid v\in V\}\\
    &\cup
    \{c_C\mid C\in\mathcal C^+\cup\mathcal C^-\}\\
    &\cup
    \{\varphi^+,\varphi^-,t^+,t^-\}
\end{aligned}
\]
und
\[
\begin{aligned}
    R_\Psi
    ={}&
    \{(v,\overline v),(\overline v,v)\mid v\in V\}\\
    &\cup
    \{(v,c_C)\mid C\in\mathcal C^+,\,v\in C\}\\
    &\cup
    \{(\overline v,c_C)
        \mid C\in\mathcal C^-,\,\neg v\in C\}\\
    &\cup
    \{(p,c_C)\mid C\in\mathcal C^+\}\\
    &\cup
    \{(\overline n,c_C)\mid C\in\mathcal C^-\}\\
    &\cup
    \{(c_C,\varphi^+)\mid C\in\mathcal C^+\}\\
    &\cup
    \{(c_C,\varphi^-)\mid C\in\mathcal C^-\}\\
    &\cup
    \{(\varphi^+,t^+),(t^+,\varphi^+)\}\\
    &\cup
    \{(\varphi^-,t^-),(t^-,\varphi^-)\}
\end{aligned}
\]
gegeben.

Wir betrachten die Unabhängigkeitsanfrage
\[
    A\perp_{\mathrm{st}}B\mid\emptyset,
    \qquad
    A=\{\varphi^+,\varphi^-\},
    \qquad
    B=X\cup\{p,n\}.
\]
Die Konstruktion ist polynomial, da für jede Variable, jede Klausel und
jedes Literalvorkommen nur konstant viele Argumente beziehungsweise
Angriffe erzeugt werden.

\paragraph{Bipartitheit.}
Setze
\[
\begin{aligned}
    U_0={}&
    \{v\mid v\in V\}
    \cup
    \{c_C\mid C\in\mathcal C^-\}
    \cup
    \{\varphi^+,t^-\},\\
    U_1={}&
    \{\overline v\mid v\in V\}
    \cup
    \{c_C\mid C\in\mathcal C^+\}
    \cup
    \{\varphi^-,t^+\}.
\end{aligned}
\]
Die Mengen \(U_0\) und \(U_1\) bilden eine Partition von
\(\mathcal A_\Psi\). Die Angriffe zwischen \(v\) und \(\overline v\)
verlaufen zwischen \(U_0\) und \(U_1\). Positive Literalargumente sowie
\(p\) liegen in \(U_0\) und greifen positive Klauselargumente in \(U_1\)
an. Negative Literalargumente sowie \(\overline n\) liegen in \(U_1\) und
greifen negative Klauselargumente in \(U_0\) an. Ferner liegen
\[
    c_C\in U_1,\quad \varphi^+\in U_0
    \qquad\text{für }C\in\mathcal C^+
\]
und
\[
    c_C\in U_0,\quad \varphi^-\in U_1
    \qquad\text{für }C\in\mathcal C^-.
\]
Auch die Angriffe
\(\varphi^+\leftrightarrow t^+\) und
\(\varphi^-\leftrightarrow t^-\) verlaufen zwischen den beiden Mengen.
Somit gilt
\[
    R_\Psi
    \subseteq
    (U_0\times U_1)\cup(U_1\times U_0),
\]
also ist \(F_\Psi\) bipartit.

\paragraph{Struktur der stabilen Labellings.}
In jedem stabilen Labelling ist für jedes \(v\in V\) genau eines der
Argumente \(v,\overline v\) mit \(\mathit{in}\) gelabelt. Wegen ihres
gegenseitigen Angriffs können nicht beide \(\mathit{in}\) sein. Wären
beide \(\mathit{out}\), hätte keines von ihnen einen
\(\mathit{in}\)-gelabelten Angreifer. Die Literalpaare kodieren daher
eine totale Wahrheitsbelegung, wobei \(v\) genau dann wahr ist, wenn das
Argument \(v\) mit \(\mathit{in}\) gelabelt ist.

Für eine positive Klausel \(C\in\mathcal C^+\) gilt
\[
    \lambda(c_C)=\mathit{out}
    \quad\Longleftrightarrow\quad
    \lambda(p)=\mathit{in}
    \ \text{oder}\
    \lambda(v)=\mathit{in}
    \text{ für ein }v\in C.
\]
Für eine negative Klausel \(C\in\mathcal C^-\) gilt entsprechend
\[
    \lambda(c_C)=\mathit{out}
    \quad\Longleftrightarrow\quad
    \lambda(\overline n)=\mathit{in}
    \ \text{oder}\
    \lambda(\overline v)=\mathit{in}
    \text{ für ein }\neg v\in C.
\]
Dies folgt jeweils daraus, dass ein Argument in einem stabilen Labelling
genau dann \(\mathit{out}\) ist, wenn es einen
\(\mathit{in}\)-gelabelten Angreifer besitzt.

Sind alle positiven Klauselargumente \(\mathit{out}\), so kann das Paar
\((\varphi^+,t^+)\) wahlweise als
\[
    (\mathit{in},\mathit{out})
    \qquad\text{oder}\qquad
    (\mathit{out},\mathit{in})
\]
gelabelt werden. Entsprechendes gilt für
\((\varphi^-,t^-)\), wenn alle negativen Klauselargumente
\(\mathit{out}\) sind.

Insbesondere lässt sich jede Belegung der Variablen aus \(V\) zu einem
stabilen Labelling erweitern. Dazu werden zunächst die Literalpaare
entsprechend der Belegung gelabelt. Anschließend wird jedes
Klauselargument genau dann mit \(\mathit{out}\) gelabelt, wenn es einen
\(\mathit{in}\)-gelabelten Angreifer besitzt. Schließlich werden
\(\varphi^+,\varphi^-\) mit \(\mathit{out}\) und
\(t^+,t^-\) mit \(\mathit{in}\) gelabelt.

\paragraph{Korrektheit.}
Angenommen, \(\Psi\) ist gültig. Seien
\(\lambda_A\) und \(\lambda_B\) beliebige stabile Labellings von
\(F_\Psi\). Wir zeigen, dass ihre Einschränkungen auf \(A\) beziehungsweise
\(B\) kombiniert werden können.

Die Labels von \(\lambda_B\) auf \(X\) bestimmen eine Belegung von \(X\).
Da \(\Psi\) gültig ist, existiert eine Belegung von \(Y\), welche unter
dieser \(X\)-Belegung alle Klauseln von \(\varphi\) erfüllt. Wir labeln
die Literalpaare aus \(X\cup Y\) entsprechend dieser Belegung und
übernehmen die Labels von \(p\) und \(n\) aus \(\lambda_B\).

Da jede positive Klausel ein wahres positives Literal und jede negative
Klausel ein wahres negatives Literal enthält, wird jedes Klauselargument
bereits von einem \(\mathit{in}\)-gelabelten ursprünglichen
Literalargument angegriffen. Dies gilt unabhängig von den übernommenen
Labels von \(p\) und \(n\). Daher sind sämtliche Klauselargumente
\(\mathit{out}\).

Nun übernehmen wir die Labels von \(\varphi^+\) und \(\varphi^-\) aus
\(\lambda_A\) und labeln \(t^+\) beziehungsweise \(t^-\) jeweils
entgegengesetzt. Nach den obigen Beobachtungen ist das erhaltene Labelling
stabil. Es stimmt auf \(A\) mit \(\lambda_A\) und auf \(B\) mit
\(\lambda_B\) überein. Folglich gilt
\[
    A\perp_{\mathrm{st}}B\mid\emptyset.
\]

Sei umgekehrt
\[
    A\perp_{\mathrm{st}}B\mid\emptyset
\]
erfüllt. Wir zeigen, dass \(\Psi\) gültig ist. Sei dazu
\(\alpha\) eine beliebige Belegung von \(X\).

Zunächst konstruieren wir ein stabiles Labelling \(\lambda_A\) mit
\[
    \lambda_A(\varphi^+)
    =
    \lambda_A(\varphi^-)
    =
    \mathit{in}.
\]
Hierzu wählen wir
\[
    \lambda_A(p)=\mathit{in},
    \qquad
    \lambda_A(n)=\mathit{out}.
\]
Dann ist auch \(\overline n\) mit \(\mathit{in}\) gelabelt. Somit greift
\(p\) alle positiven und \(\overline n\) alle negativen
Klauselargumente an. Sämtliche Klauselargumente können daher mit
\(\mathit{out}\), die Argumente \(\varphi^+,\varphi^-\) mit
\(\mathit{in}\) und die Argumente \(t^+,t^-\) mit
\(\mathit{out}\) gelabelt werden. Die übrigen Variablenpaare können
beliebig konsistent gelabelt werden. Das resultierende Labelling ist
stabil.

Weiter wählen wir ein stabiles Labelling \(\lambda_B\), das auf \(X\)
die Belegung \(\alpha\) kodiert und
\[
    \lambda_B(p)=\mathit{out},
    \qquad
    \lambda_B(n)=\mathit{in}
\]
erfüllt. Ein solches Labelling existiert nach der obigen
Vervollständigungsbeobachtung: Die Variablen aus \(Y\) können beliebig
belegt und die übrigen Argumente anschließend entsprechend gelabelt
werden.

Aus
\(A\perp_{\mathrm{st}}B\mid\emptyset\) folgt nun die Existenz eines
stabilen Labellings \(\lambda\) mit
\[
    \lambda|_A=\lambda_A|_A
    \qquad\text{und}\qquad
    \lambda|_B=\lambda_B|_B.
\]
Insbesondere gilt
\[
    \lambda(\varphi^+)
    =
    \lambda(\varphi^-)
    =
    \mathit{in},
    \qquad
    \lambda(p)=\mathit{out},
    \qquad
    \lambda(n)=\mathit{in}.
\]
Da \(\varphi^+\) und \(\varphi^-\) mit \(\mathit{in}\) gelabelt sind,
müssen alle positiven und negativen Klauselargumente
\(\mathit{out}\) sein.

Sei \(C\in\mathcal C^+\). Da \(c_C\) mit \(\mathit{out}\) gelabelt ist,
muss es einen \(\mathit{in}\)-gelabelten Angreifer besitzen. Wegen
\(\lambda(p)=\mathit{out}\) muss dies ein Argument \(v\) mit \(v\in C\)
sein. Somit ist \(C\) unter der durch \(\lambda\) kodierten Belegung
erfüllt.

Sei nun \(C\in\mathcal C^-\). Aus
\(\lambda(n)=\mathit{in}\) folgt
\(\lambda(\overline n)=\mathit{out}\). Da \(c_C\) mit
\(\mathit{out}\) gelabelt ist, muss daher ein Argument
\(\overline v\) mit \(\neg v\in C\) mit \(\mathit{in}\) gelabelt sein.
Damit ist auch \(C\) erfüllt.

Die von \(\lambda\) kodierte Belegung von \(Y\) erfüllt folglich alle
Klauseln von \(\varphi\). Ihre Einschränkung auf \(X\) ist
\(\alpha\). Da \(\alpha\) beliebig gewählt war, gilt
\[
    \forall X\exists Y\,\varphi.
\]
Somit ist \(\Psi\) gültig.

Wir haben damit gezeigt, dass
\[
    \Psi\text{ gültig}
    \quad\Longleftrightarrow\quad
    A\perp_{\mathrm{st}}B\mid\emptyset
    \text{ in }F_\Psi.
\]
Da die Konstruktion polynomial und \(F_\Psi\) bipartit ist, folgt die
\(\Pi_2^{\mathrm P}\)-Härte für bipartite AFs. Zusammen mit der
Zugehörigkeit zu \(\Pi_2^{\mathrm P}\) ergibt sich die Behauptung.
\end{proof}